\section{The Quality Known as Repose}

\begin{quotationx}
The quality known as \emph{repose} in the ancient Greek is a manifestation of that serenity which belongs to a people not yet disturbed by self-doubt, self-immolation and self-contempt. It is the extreme harmony of a mentality not yet shaken by the tortures of introspection or inner conflict, by what Goethe called `two souls throning within one bosom'. The beauty of the Greeks is the beauty of men who have never in their wildest dreams beheld the horrors of Dante's Inferno.

Poorer than the moderns in this respect, they consequently have the bliss which is partly ignorance, and this bliss is revealed in their art. Everything that has appeared in western Europe since the fall of the Roman Empire is certainly less serene, less blissful, more foolish, perhaps, in its wisdom, than was the partial ignorance of the Greeks; but it is more fretful, more nervous, more subterranean and subcutanesous, more full of insight and second sight, and consequently, therefore, more disturbing.
\flright{\textsc{Anthony Ludovici}, \emph{Personal Reminiscences of Auguste Rodin}}
\end{quotationx}



\flrightit{Posted on 2010-06-19 by Anonymous}

\begin{center}* * *\end{center}

\begin{footnotesize}
\begin{sffamily}
\texttt{kadambari on 2010-06-20 at 09:50 said:}

``than was the partial ignorance of the Greeks''

Partial ignorance of the Greeks? This is seems a Christian prejudice. Muslims also keep saying that Islam is an improvement over the pagan civilizations--perhaps for Arabs but not for other on whom this religion was imposed. The Greeks did not know the self-loathing that arises from a conception of God as ``exteror'' and they did not know the ``revealed truths'' of the semitic faiths that results in guilt and self-loathing. One of my Iranian friends told me how after Islam everything beautiful seemed to disappear from Persia--beauty, song, dance, art (how Islam reminds her of death with its emphasis on covering women in black and distaste for art) and how the Persians kept their identity by not losing their language despite the conversion. It must be the same for Europe. Only someone who does not understand the civilization of the Greeks could think of the fires of hell of Christianity as an improvement. One thinks of what would have happened rather if the semetic faiths never came into being and disrupted civilizations in between. The East and West would have interacted harmoniously. Look at the harmony of the short lived Hellenic-Buddhist civilization in Afghanistan for a while which was then completely destroyed by barbarians. The two cultures interacted harmoniously. The contemplative civilization of the Greeks was destroyed--even the Romans were quite a different peoples--only when Greek ideals resurfaced (even by influencing Christianity) has the Western world created greatness... .Why is it a strength to have self-doubt, self-immolation and self-contempt? Someone please explain? I am a thousand times more enlightened by reading Plato and Aristotle than the Bible (everything that is good in it can already be found in the the older religions such as Hinduism and Buddhism.) I do not see why the route to spirituality has to be via the Christian one and how everything that came earlier needs to be ``partial-truths''? However, I realize that Europe has been Christian and that its civilization has been Christian. But to dismiss what came earlier as ``lacking'' seems to be to lacking in understanding of what came before ``Christianity''. However, as one of my Persian friends told me of how she thinks of Islam as something not Persian, but since they have been converted and have lost their old religion, they have no other religion to call their own so they can't deny this religion as they are left with nothing. I sense that it must be the same way for Europe--especially as the conversion of Europe was total. There was definitely a loss there--which only those pagans still left in the East can sense as they have not lost the pagan heritage... .

\hfill

\texttt{Cologero on 2010-06-20 at 12:03 said:}

Please be cautious in your reading. Often, a commitment to a certain ideological perspective inhibits one's ability to understand certain passages. The ignorance in question refers to the ``tortures of introspection or inner conflict.'' Ludovici is not commending that lack of ``repose'' that afflicts the modern.

And just to point out that life and thought are not as simple as ideology assumes, Ludovici claims that ``Christianity is merely Platonism for the mob.'' As a matter of fact, he blames Socrates for the disappearance of beauty in the West, because he exalted the soul and the invisible over the body. In this regard, Ludovici is following a line of thought from Nietzsche.

\hfill

\texttt{kadambari on 2010-06-20 at 12:24 said:}

they consequently have the bliss which is partly ignorance

You did not answer my question but answer that I hold an ideology. I do not. How is their bliss ``ignorance''? Is it not a prejudice on his part to call their art stemming from ``ignorance''? I do not have any ideological perspectives. I have merely noticed that it is often common for Muslims and Christains to view their world view as somehow an ``improvement'' over the pagan one. And of course, the rigorous Christian is often on a higher level than the new age liberal without a doubt.

\hfill

\texttt{Cologero on 2010-06-20 at 22:40 said:}

I certainly did answer your question about ``ignorance''. You force me to repeat myself: \emph{The ignorance in question refers to the ``tortures of introspection or inner conflict.''}

The other point I made, which again you ignore, is that that ignorance which Ludovici refers to has no relationship at all to the alleged superiority of ``Christianity''. Since that conclusion can be found nowhere in the text itself, then it has to have come from your own presuppositions.

\hfill

\texttt{kadambari on 2010-06-21 at 22:55 said:}

How can you be so sure that the bliss on their faces is not something superior?

`The beauty of the Greeks is the beauty of men who have never in their wildest dreams beheld the horrors of Dante's Inferno.'

Again why does introspection have to arise from external fears?

Actually the repose reminds me of the serenity you see on the faces of Hindu deities... What if it is trying to express a serenity alien to the Christian mindset?


\hfill

\end{sffamily}
\end{footnotesize}

